\chapter{亡故者(仍未死去……)}

\section{事件——或身体与精神}

\notehead{98}{}

\dmark{}(9月22日)葬礼各注中日期最新的一条(几条脚注不算),是5月24日——距今已有四个月。其后的两个星期,直到6月10日,主要用于重读并补充、修改已写成的各注;其间还有Zoghman~Mebkhout一两天的来访:他来通读葬礼的全部各注,在我把它们交付打印之前,并给我他的评点。\origpage{421}我原以为,定稿会在6月初备齐,并(这还是乐观的……)在大学暑假大忙之前打好印出。我满心想在人人打点行装之前,把我的“五百页的信”寄给各位!

实际上,在我写这些文字的时刻,葬礼的文本仍未完成:一如四个月前,它还缺最后的两三条注——外加其间新添的一条\fnote{1}{}{(9月23日)实际上,看来这条预定要写的“注”,已爆成三条不同的注(n\textsuperscript{o}~99--101)。}:即我方才以下列文字着手的那一条,算是这段时间里所发生之事的一份简报。6月10日,一桩新的意外闯入《收获与播种》——这部从不缺意外的写作:我病倒了!一阵胁痛突然袭来(前一分钟我毫无预感),以不容分诉的力量把我按倒在床上。站立乃至坐姿,顿时变得十分痛苦;唯有躺卧似还可忍。真是蠢透了——尤其在此刻:一项十万火急的工作眼看就要收尾,更不必说别的!躺卧着打字,没门;就连躺着写字,也不是件轻松事……

我又花了将近两个星期——其间我竭力顶住一切、照旧工作——才向这明摆着的事实低头:我的身体已经衰竭,它执意要求——而我不肯听闻——完全的休息。

我之所以这么难以听话入耳,是因为我的精神依然清新、敏捷,跃跃欲试要继续冲刺,仿佛它有一条独立于身体的自主生命。它甚至清新敏捷到这种地步:连身体的睡眠需要都难以兼顾,一再把睡眠——这个搅局的煞星!——的期限,推到衰竭的极限。

我一生中,直到三四年前,经由深长睡眠而有的无限的恢复力,曾是那些有时失度的精力投资的坚实而救命的配重:只要睡眠有保证,人就无所畏惧;尽可以(而不失为清醒)全身心投入工作的狂欢、直到衰竭——大不了再用修复性睡眠的狂欢补回来!这份能力——我一生把它与工作能力、与发现的能力(两者想必血肉相连……)一并视为理所当然——近年终于磨损,有时甚至消失;缘由我此刻仍辨不真切,也还不曾真正下力气去探它。越来越常见的是:在一整天的打字机(或手写笔记)之后,顺从我那拒绝继续的身体之命,我决意去睡;可一躺下(连同躺卧对坐姿紧张的局部缓解),立刻重新点燃反思。它愈发精神,一连几个小时,乃至整夜(或者不如说,夜里所剩的那点时间……)。我尽管明白这套系统无利可图(姑且不说它长远是否“活得下去”),——因为(至少在我这里)没有书写支撑的长久反思,终归原地打转,常沦为一种反刍——恶习已然养成,且有加剧之势。我觉得,它成了我近些年生活中能量涣散的大本营;而别的涣散机制,却经年累月一一被清除。

倘若这一机制在我生活中扎下如此顽固的根,倘若这些年我一直肯付这样的代价,那肯定是我身上的某种东西从中得到了好处,时候一到还会再从中获益。凑近些检视这一局面,并非奢侈——四个月来,我不止一次险些动手。

\dmark{}这想必是一件紧迫的事。然而我终于懂得:还有更紧迫的。\origpage{423}我得先应付燃眉之急:重新接上与我身体断掉的接触,帮它从我终于感到并承认的衰竭状态里爬出来,找回失掉的元气。我那时懂得:为此,我必须在一段不确定的时间里,放弃一切智力活动——哪怕是沉思临我之事的意味。今天重新提笔的这些注,就是为它画上句号:我几番大投资中这段漫长而有益的“括号”——那些投资,曾有一段(自今年2月起)汇入“收获与播种”的写作。本注是对这四个月“括号”的最初一点思考,或无论如何,一份简要的汇报。

等到终于懂得需要完全的休息,大的疲惫已成了深深的衰竭。由于不曾肯听我身体那分明不容分诉的语言,最初两个星期里,我在体力疲惫的状态下硬挤出来的那几页可怜的评注与修改,是以一份回头看来近乎疯狂的精力投入换来的!总之,在这些壮举之后,我只得长卧数星期,每天只起身几小时,做些必不可少的实际事务。

引人注目的是,一旦终于\textbf{懂得}了完全休息的需要,我便毫不费力地彻底脱离了一切智力活动,没有丝毫“作弊”的念头。我甚至无须真正做什么决定——单是懂得,我就已然脱开了。前一天还吊着我胃口的那些任务,忽然显得十分遥远,仿佛属于一段非常遥远的过去……

然而当下并非空无。几个星期乃至几个月里,睡眠迟迟不肯来;我长卧数小时,表面全然无所事事,却不记得有哪一次觉得时间难熬。我与我的身体重新相识,也与最切近的环境重新相识——我的房间;有时是眼前那片晒着太阳的草皮或干草:我在屋旁、或在一次短暂(而谨慎……)的散步中随意躺下的地方。我久久地追随一只苍蝇在一束阳光里的舞蹈,或一只蚂蚁、或一些绿色或粉红色的半透明小虫在无穷草茎上的长征——那些草茎在我眼前缠成难解的丛林\dmark{}\origpage{424}。也正是这样的心境:借着寂静与大疲乏之便,人以关切之心,追随最轻微的风在自家肠肚里的踌躇巡游——总而言之,是这样的心境:人们借此与最基本、最本质的事物重新接触;在这样的心境里,人们充分掂量一场修复性睡眠的全部益处,甚至掂量“能够无碍地撒一泡尿”竟是怎样的奇迹!

身体的寒微运转是一件非凡的奇迹;只有当这运转以这样或那样的方式受到扰乱时,人们(有时违背本心)才会多少意识到它。很清楚,“技术上”说,我“健康问题”的底子,是睡眠之扰。这一扰动的深层缘由,当时向我滑脱,如今依然。我是摸索着,首先试图找回睡眠——那种我熟识的酣畅大觉;而它偏偏在我最需要它的时刻神秘地躲避我!直到不久前我才找回它。想必不必说:求助于药片的念头,根本不会来到我心头;即便我试过一些汤剂或橙花水(借这次机缘认识的),我心底也明白那只是权宜之计。更正经的一件事:我借此机会对饮食做了重要的调整:减少淀粉类,增加绿叶蔬菜与水果(生熟皆可);(适度地)重新把肉作为日常饮食的常客;尤其是,极其严格地削减脂肪与糖的消费——在这方面,我(一如富裕国家里的许多人)至少自战后以来,一直有一种系统性的失衡。帮助我认识到这种饮食调整对找回失衡生活之平衡的重要性——尤其还有每日几小时量级的体力活动的必要,以抗衡紧张的智力活动;否则后者动辄耗竭身体,把可用的生命力往上抽,造成强的阳失衡——的,是我的女婿艾哈迈德;他行中医,对这类事情有很好的“感觉”。

况且,艾哈迈德并不满足于给我一些好建议,外加一套阴阳辩证——四五年来,我既有充分的机会熟悉事物这份精微的动力学,对它相当敏感。一旦我身体好转到可以打理园子,\dmark{}\origpage{425}而他看到我亲力亲为、想把一处寒碜的小园子多少重新开动起来,他便抢先发起几项大面积的工程:开出新的地块,运土,移栽与播种,做梯田、垒挡土矮墙,重整堆肥……随着日子与星期流逝,我眼看这些整治任务在我这位不知疲倦的朋友的推动下于眼前铺展——它们足以让我忙上几年,乃至忙完余生!

这恰是我所需要的,也是长远而言我所需要的,用以配衡一种过于炽烈的智力活动。在这方面,每天强迫自己散步——如人们早就这样建议过我——帮助不大:散步时脑袋照旧研磨,一如在床上,几乎视而不见地穿过景色的美!相反,浇园——我得当心园子的长势;更好是锄一畦菜——我就不由自主地留心、并多少上心:体会土的质地,锄地如何改变它,菜蔬与“莠”草如何影响它,堆肥与铺盖如何影响它;久而久之,也体会我理应照料的那些植物的境况——那境况,在很大程度上,映照着我肯给它们多少留意。园艺这项活动,以及环绕它的一切,回应我身上两股强烈的渴望或心性:其一推我去从事一种每天能看见\textbf{手底下出东西}的行动(散步绝不是这样;某位同行兼朋友建议的举重哑铃,更不是……);其二推我去从事一种每时每刻都有机会\textbf{学习}的行动——在与事物的接触中学习。我似乎恰恰在“做”某件事情的境况里最善于学习——某个在我手底下成形、变形的“某件事情”……

一旦跨过严格意义的衰竭状态,我的康复——依我看——得益于两类活动;或不如说,得益于日常活动(屋里的与园里的)中两类重要而有益的因素。其一是\textbf{体力的付出}:哪怕动工之前我常觉疲惫、没劲——活越是“重”,叫我抡起\dmark{}\origpage{426}重镐或搬动大石头的——过后我越觉精神,带着一份沉沉的好疲惫。其二是与\textbf{活物}的接触:需要照料的植物;需要为它们备好、随后铺盖或锄松的土地;需要烹备、而我吃着与备餐时同样有滋味的食物;讨食、也讨它那份抚爱的猫;还有各样器皿与工具,乃至那些磕磕碰碰、常不甚规整的石块——得把它们翻来转去,垒成肯站得住的矮墙……

体力的付出与活物的接触——这恰是智力工作所缺的两个方面;正因如此,智力工作按其本性是不完整的、残缺的,走极端时——若不以别的东西补足与抵偿——甚至危险乃至有害。这是短短三年多来我第三回有机会认识到这一点。如今甚至十分清楚:我面对一道苛刻的期限:改变某种生存方式,找回一种平衡——使我存在的阴极、我的身体,不再被持续怠慢,而一味让给阳极、精神(或更直白地说,脑袋)——否则,就在最近几年里把命搭上。这是我的身体对我说的,说得再清楚不过!我如今到了生命的一个节点:某种基本“智慧”的需要,已成了一个\textbf{活命}的问题——按这个词的本来与字面的意义。这肯定是好事——否则,所谓“智慧”将永远被推给“下回”,让位给那种智力活动上的贪食症;那种贪食,是我整个成年生活中占支配地位的力量之一。

面对这样一道分明的期限——“改变,否则就死!”——我无须自问选择为何。正因如此,将近四个月里,我得以在从不觉得勉强自己的情况下,戒绝一切智力活动——数学也罢,非数学也罢。我没有对自己说,却知道:说到底,一个活着的园丁,总比一个死去的数学家(或一个死去的“哲学家”或“作家”,反正一样!)值钱。带点促狭,还可以补一句:甚至比一个活着的数学家还值钱!(不过,那就是另一段故事了……)

况且,我不认为哪天自己会被逼进这样一种“极限”境地:长远放弃一切智力活动——数学的也罢,沉思的也罢。毋宁说,未来这些年里最直接、最紧迫的实际任务,依我看恰恰是达到一种生活的平衡:两类活动——身体的与精神的——日复一日共存,谁也不变得吞噬一切、排挤对方。\origpage{427}我不自欺:自我童年起,我最强的投资正在“精神”这一侧;今天仍支配我生活的两大热情,也仍把我引向那一侧。这两大热情——数学的热情与禅观的热情——依我看,尤其是前者,若非唯一地,充当了我生活中的失衡因子:像某种至今仍带着那份恼人的“吞噬其余、独占一切”的倾向的东西。肯定并非偶然:1981年6月以来,我生活中标志着失衡局面的三桩“患病情节”,恰恰都落在数学热情当台主演的时期。

人们会说,最后这桩情节——它发生于《收获与播种》的写作之中——并不完全如此:那部书构成一场对我自己的反思,甚至不失为一场名副其实的禅观。但同样真确的是:这场对我数学家过去的反思,一直由我的数学热情喂养;尤其在那第二部分——葬礼——依我看,这份热情的自我本位成分,以格外强而持续的方式被牵扯在内。然而,即便回顾,我并不觉得这场反思曾在任何时刻采取一种吞噬性的、乃至疯狂的节奏与音高,一如前两次——那时我的身体终究被迫喊出一声不容分诉的“受够了!”。与我的一生分割开来看,我一年半以来的智力活动(自“重启”以来——先《追寻栈》的写作,继以《收获与播种》)显得以相当合理的节奏进行,没有忘了喝水,也没有忘了吃饭(不过有时,稍稍,忘了睡……)。倘若它终于酿成第三桩“健康情节”(用个委婉语),那无疑是植根于整整一生:那永恒的失衡——一颗过强的脑袋,把它的节奏与法则强加给一个壮健的、长期默默承受的身体\fnote{2}{(*)}{这里我应当把1974至1978这五年算作例外:那几年没有被什么大任务主宰,手工的活计占去了我时间与精力的不小一部分。}。

\dmark{}\origpage{428}过去两个月里,我充分体会了身体劳动那无可替代的益处:与朴素的活物亲密接触,听它们默默讲述书本或单凭反思无力教给我的简单而根本的事。借着这样的劳动,我重新找回了睡眠——这位比吃喝还要珍贵的伙伴;也随之恢复了元气,重获那一度仿佛突然消失的健壮。我还发觉,在我目前的人生阶段,若想把去年开始的这场新的数学冒险再继续几年,就非得把双脚稳稳扎在园子的沃土里,否则便会危及我的健康与生命。

未来几个月,我得建立一种新的生活方式,使身体的劳动与精神的劳动每天都有一席之地,彼此调和。要做的事多着呢!

\section{陷阱——或省力与衰竭}

\notehead{99}{}

(9月23日)昨晚我不得不强行收住,免得顺着势头写到凌晨两三点,重新落回那副我熟得不能再熟的齿轮。我神清气爽;倘若顺着自然的坡道,我甚至会写到拂晓!智力工作的陷阱——至少那种怀着热情、在一门久经熏陶终于如鱼得水的材料里从事的智力工作——就在于它\textbf{容易}得令人难以置信。拉吧,拉吧,东西总会来;只消拉就是;有时才勉强感到一点吃力、一点摩擦,表明它多少有点抵抗……

然而我记得,在我做数学家的年轻岁月,有一种挥之不去的沉重感、滞重感,须以顽强的努力去克服,努力过后留下一阵疲惫。那主要对应于我一生中工具不足、乃至不顺手的时期;或对应于稍后的另一段:当时我不得不多少痛苦地、在一种(主要是Bourbaki小组的)环境的压力下,习得一套多少“四面八方”的工具——那环境熟练地使用它们,而它们的所以然却不曾随时显现给我,有时甚至多年如此。我已有机会在《收获与播种》第一部分里,谈到那些有时颇为费力的年头(见“受欢迎的外来者”第9节,与“炉中百铁,或:干熬无济于事!”一注,n\textsuperscript{o}~10)。那主要是1945至1955年,与我的泛函分析时期重合\dmark{}\origpage{429}。(我觉得,在我1960至1970年间的学生们身上,这种对缺乏足够动机的学习——凭长辈的权威硬吞观念与技术——的抵抗,比我当年弱得多;干脆说,我全然没有觉察到。)

回到本题:尤其自1955年起,我常常有“飞翔”之感——做着数学,毫不费力,一如我昔日曾那样艳羡过的某些长辈,艳羡他们那份近乎奇迹的从容——那份从容,对我这个寒微而滞重的身量来说,显得全然高不可攀!今天我明白:这样的“从容”并非某种特殊天赋的特权(这样的“天赋”,我曾在某些人身上、在我自己毫无此种“天赋”的时刻见过);它宁可说是自行显现为一份果实:一份对某门材料(比如说数学)的炽热兴趣,与对其或多或少的长期熏陶之结合的果实。倘若“天赋”确在这份从容的出现中插了一手,那想必是经由时间这个因子——或长或短,因人而异(有时在同一个人身上,也因时机而异,这是真的……)——才得以在某一题目的工作中抵达完美的从容\fnote{3}{(*)}{然而,我认识好几位数学家,各自做出过深刻的工作,却从不曾给我这里所说的那份从容、“轻便”的印象——他们仿佛时时与一种无所不在的滞重搏斗,每走一步都得费力克服。由于这种或那种缘由,方才所说的那份“自然的果实”,在这些杰出的人身上不曾“自行显现”,尽管它本该如此。可见并非所有的结合都结果实……}。

总之,年岁越往前,我做数学时这份“轻便”之感越强——仿佛事物只消我们肯费心看一看、稍稍端详,便只等着向我们显露自己。这不是技术精湛与否的问题——分明,从这个角度看,我的状态比1970年“退出数学”时差得多:此后我有机会做的,主要是“忘却”当年所学;只在私下里偶一为之“做做数学”,而且怀着一种与往昔(至少乍看)十分不同的精神、就着十分不同的主题。我也不是说:只消我肯去啃某道名题(比如Fermat的、Riemann的或Poincaré的),就能在一两年\dmark{}\origpage{430}乃至三年之内\fnote{4}{(*)}{[原书此注文字与注3完全相同,照录。]然而,我认识好几位数学家,各自做出过深刻的工作,却从不曾给我这里所说的那份从容、“轻便”的印象——他们仿佛时时与一种无所不在的滞重搏斗,每走一步都得费力克服。由于这种或那种缘由,方才所说的那份“自然的果实”,在这些杰出的人身上不曾“自行显现”,尽管它本该如此。可见并非所有的结合都结果实……},笔直地开辟一条通往其解答的路!我所说的那份从容,不是那种通往一个预先定好的\textbf{目标}——证明某条猜想或给它一个反例……——的从容;它宁可说是那种使人得以跃入未知的从容:朝着某个方向——一个幽暗的本能告诉我们它是丰产的;带着一份亲身的把握——它从不会被辜负:旅途中的每一天每一时,都不会不带来一份新的认识的收成。\textbf{怎样的}认识恰好等着明天,乃至今天的下一个钟头,我们固然有所预感;而正是这份不断被赶到前头的“预感”、与它浑然一体的悬念,不断把我们抛向前方——此时,我们所开掘的事物本身,仿佛在把我们吸进去。已然知道的,总在精确、滋味与丰富上超出所预感的;而这份已知,又旋即成为起点与材料,供一份焕新的预感腾身而起,追逐一个渴望被认识的新未知。在这场发现事物的游戏里,我们每一刻所循的\textbf{方向}是已知的,而\textbf{目标}却被忘记——假如我们确实曾从一个目标出发、曾打算抵达它的话。这个“目标”,其实当时是一个\textbf{出发点},一份野心或无知的产物;它尽过自己的职责:推动“老板”,定一个初始方向,并引发这场游戏——在这场游戏里,目标其实没有份。只要踏上旅途不是为了一两天,而是长程的,那么时日与月份的流逝将向我们揭示什么、在一长串未知的波折尽头将把我们引向何处,对旅人是一个全然的谜;一个如此遥远、老实说如此可望而不可即的谜,以致他不大操心!倘若他有时凝望地平线,不是为了那件不可能的差事——预测一个终点,更不是为了随心所欲地裁定它——而是为了定位他此刻正到哪里,并在供他继续旅程的诸方向之中,选出从今往后他感到最灼人的那一个……

这就是我方才谈及的“难以置信的轻便”——就一门全然智性的工作(如数学)中的发现而言。它既不被\textbf{内在的抵抗}\fnote{5}{(*)}{然而,我认识一位天资出众的数学家,他与数学的关系是典型的冲突型的:每走一步都被强大的抵抗所阻滞,比如害怕某一期望(比如以一条猜想的形态)到头来被证明是假的。这样的抵抗,有时可以酿成一种真正的智力瘫痪。可与前一条脚注对读。}所\textbf{阻滞}(一如它在我所实践的禅观工作中如此常见),也不要求付出一种\textbf{体力}——那体力会产生疲劳,而疲劳终于发出毫不含糊的停止信号\dmark{}\origpage{431}。至于\textbf{智力}的付出(假定到了只剩“时间因子”这一层“抵抗”的地步,还能谈“付出”的话……),它似乎既不产生智力的疲劳,也不产生身体的疲劳。更确切地说:即便有身体的“疲劳”,人们也并不真把它当作疲劳来感受——除非偶尔的腰酸背痛(因长时间保持一个固定坐姿),以及诸如此类的附带麻烦。这些,换个姿势便轻易消除。而躺卧的姿势有一种不合时宜的功效:叫它们烟消云散,从而促成智力工作的重新开工——而不是那份亟需的睡眠!

然而,我终于明白:有一种比肌肉或神经的疲劳更精微、更阴险的身体“疲劳”——后者以无可辩驳的休息与睡眠之需,显出自身。这里的字眼,与其用“疲惫”,不如用“衰竭”更切;当然,要说定:这一状态并不按这个词的通常意义被知觉——通常意义指极端的疲惫,表现为光是起身、走几步路都须巨大的努力。这里说的是身体的能量为大脑所付出的一种“衰竭”;它表现为身体总的“张力”、其生命力水平的逐步低落。看来,由过度的智力活动(我是指:不以足够的身体活动——那种产生身体疲劳与休息之需的活动——来补偿的智力活动)造成的这种衰竭,是渐进而\textbf{累积}的。这些效果,同时取决于给定时期内智力活动的\textbf{强度}与\textbf{时长}。以我从事智力活动的强度,以我的年纪与体质,这一累积衰竭,似乎在一两年不间断的活动之后——没有规律的身体活动作补偿——就抵达一个危险的临界阈值。

在某种意义上,我所说的这份“轻便”只是表象。紧张的智力活动动用可观的能量,这很清楚:能量在某处被取用、“花费”在一项工作上。这个“某处”,似乎就在身体的层面:身体竭其所能“照单收下”(或不如说,\textbf{垫付})脑袋那不计成本地自开的(有时令人眩晕的)开销。身体所供能量正常的回收途径,是睡眠。脑袋一旦变得贪食,终于就要侵占睡眠——这就等于吞噬一份不再更新的能量资本。\dmark{}\origpage{432}智力工作之“轻便”的陷阱与危险在于:它孜孜不倦地怂恿我们越过这道阈值,或越过之后停留在外;而且这一越界,并不以疲劳、乃至衰竭的那些惯常而确凿的征兆,向我们报警。我明白:要察觉这道阈值的临近与越过,需要极大的警醒——而那时,人正全身心投入一场迷人的冒险。在身体的层面知觉这份能量的亏空,要求一种对身体的倾听状态;我常常欠缺它,也很少有人具备它。况且我怀疑:这样一种有意识的注意与身体之间的契合状态,能在任何人身上、在其生命中一段纯然智性活动独占、排斥一切身体活动的时期里绽放。

况且,许多智力工作者凭本能感到这种身体活动的需要,并相应安排生活:园子、木工、山、船、运动……像我这样,为了一个过于侵占的热情(或一种过强的倦怠)而忽视这份健全本能的人,迟早要付账。三年里,这是我第三回结账;我必须说,我付得并无怨言——或者更确切地说,心怀感激:每一桩新的患病情节都让我明白,我不过是在收割自己疏忽的果实;而且,它还带给我一份教益——那份教益,恐怕唯有它能给我。这最后一桩——它刚刚结束——所给我的主要教诲,也许是:是时候抢在前面,叫这样的“提点”从此不再必要;或者更具体地说:是时候照料我的园子了!

\section{向Claude~Chevalley告别}

\notehead{100}{[原书此处误印作“注~97”]}

昨天与今天的反思里,我起初有意搁下了一桩事件:它正落在患病情节的正中,在7月的最初几天——也就是说,当时我还卧床。那就是Claude~Chevalley的逝世。

我是经由《解放报》上一篇多少以此为题的含糊文章得知的;一位女友随手递给我,想着也许我会感兴趣。关于Chevalley,文中几乎什么都没有;倒有关于Bourbaki的几大段——他是其创始成员之一。得知消息时,我感到自己全然发懵。几个月来,我总看见自己眼看就要做完《收获与播种》——打好、印出、装订,一应俱全——然后火速上\dmark{}\origpage{433}巴黎,把一份还热乎的副本送到他手里!倘若天底下有一个人,我确信会怀着真正的兴趣、且常常怀着乐趣读我那块铺路石,那就是他;而且我全然没有把握,除他之外还有没有第二个!

自反思之初,我就发觉:Chevalley在我行程的一个关键时刻,给过我某种东西——某种播于一场沸腾之中、却在寂静里发芽的东西。当时我感到把我与他系在一起的那种东西,与其说是一份\textbf{感情}——比如感激、同情或爱——不如说是别样。这些感情当然在场,一如它们对那些早在二十多年前把我当作“自己人”接待我的“长辈们”中的这一位或那一位,也在场。使我与Chevalley的关系有别于我与他们中任何一位、与我的大部分朋友(不妨说与所有朋友)的关系的,是另一样东西。那是一种感觉,或者更恰切地说,一份知觉:一种\textbf{本质的亲缘}——超越文化的差异、超越自我们少年时代起便打在我们身上的种种烙印。我说不好,在我反思中谈到他的那几行文字里\fnote{6}{(*)}{见“与Claude~Chevalley相遇——或:自由与善良的情感”(第11节),及下一节“功劳与轻蔑”的最末一段。},是否透出这份“亲缘”的些许。在我生平那段——即那几行所指的——时期里,Chevalley显出的,也许更是一位“长辈”——这一次,是在对生活中某些基本事物的领会层面上——而不太是一位“亲人”。然而这份距离,我后来的成熟想必把它缩小、乃至取消了;一如在数学层面上,这早已是旧事——在我与他的关系中,一如在我与其他长辈的关系中。倘若我此刻试图以言词勾勒这份亲缘的意味,或至少它的一个征兆,来到心头的是这个:他与我,都是“独行的骑士”——各自走在自己的“孤独的历险”里的旅人。关于我自己的那一场,我在《愚妄与更新》最后那一章(同名)里有所陈述\fnote{7}{(**)}{这个方向上,尤其见“禁果”与“孤独的历险”两节,n\textsuperscript{o}~46、47.}。也许,对那些深知Chevalley的人(乃至对别人),反思的这一部分,比专门谈他的那一部分,更能暗示我想表达的东西。

与他相见、并与他稍稍谈谈,本必定使我比过去更好地把握这位朋友,更好地安放这份本质的亲缘,以及我们的差异。倘若除了Pierre~Deligne之外,还有一个人让我急于把《收获与播种》的文本亲手交到他手里,那正是Claude~Chevalley。\dmark{}\origpage{434}倘若有一个人的评点——无论是顽皮的还是辛辣的——会对我有特殊的分量,那还是他。7月的头一个星期的那一天,我知道:我不会有这份乐趣了——再不能把我最好的东西带给他,也再听不到他的声音了。

那件怪事——它想必助长了我得知消息时的那份\textbf{发懵}——是:在此前的几个月里,不止一次,在想着与Chevalley的一场临近的会面时,我记起他正受着健康的困扰;我心中有一种被不断推开的不安:那场会面也许不会有了,我的朋友也许会在我去看他之前消失。给他写信或打电话的念头当然掠过心头——哪怕只为问问他的健康、他的近况,对他说几句我正投入的那项工作、以及我打算为此去看他的事。而我把这个念头当作愚蠢、当作打扰推开了(反正真没有什么理由……等等),一如人们在诸如此类的情形里常常做的那样。这恰好例证:我这个人,一如许多别人,仍未能充分运用自己的能力——推开那股对事物的幽暗预知;它向我低语一份认识,而我太忙、也太懒,不肯去听……

\section{表面与深处}

\notehead{101}{}

(9月24日)前两天围着“患病情节”兜了两天圈子之后,该是接上6月间中断的线头的时候了——就从我搁下的地方接起。我那时预计,还有最后两条注要写:一条“葬礼悼词(2)”(它将接续并补全5月12日的“葬礼悼词(1)——或恭维”),以及一条终曲“De~Profundis”,我打算在其中为围绕葬礼的整场反思勾勒一份盘点。

这两条注预定要写的内容,在我病倒的那一刻还热乎着——我眼看就要把它们统统倒到纸上,只差把最后一手加到前面所有各注上,好有这种在“稳固而齐整的后方”上工作的感觉……在其后整整三个月里(自6月23日确切算起),我实际上停止了葬礼上的一切工作,只除偶尔的几处排印更正;它(唉)有点从我脑子里溜走了。我甚至觉得自己有点傻——总之有点难为情——竟一本正经地坐下来,去填那些在“罚抄标题”后面等着的一页页白纸,只凭这些标题出现在一份临时的目录里,而我曾轻率地在某份注定要发表的文本里此处彼处提到过它们。\dmark{}\origpage{435}尤其是“葬礼悼词(2)”的情形;就连方才重读初稿“葬礼悼词(1)”(别名“恭维”),也不曾足以焐热那份在它的角落里从容冷却了几个月的实体!

然而,就在写下那条注的5月12日的次日,及其后整整一个月里,我的手一直发痒,想更深地开掘那座我刚刚不知不觉按上手的新矿。去年,当Nico~Kuiper好意给我寄来IHES建所二十五周年的小册子时,我花了大约半个钟头翻了一遍(包括关于Deligne与我各半页的两篇短文),没发现什么特别的。唯一打动我的,是通篇没有一句影射IHES那艰难的最初年月——那时它的名声是在一间凑合的屋子里立起来的,而我(连同最早的几届代数几何研讨班)是唯一在“实地”代表它的人。几个月后,1984年3月,我写“有益的挣脱”一注(n\textsuperscript{o}~14)时,又想起了它。因为不敢相信自己的记忆,我出于求个心安,请尼科再寄我一册那本小册子(第一册已找不着了)。这是第二回翻那两篇短文,眼睛也许不那么匆忙了。然而,这一回我又没“入局”,实在是。顺带我多少惊讶地注意到:关于Deligne的那篇里说,“其工作的主导轴心是‘理解代数簇的上同调’”——谁信呢!而这事,我一忘就是一两个月(直到写“拒受遗产——或矛盾的代价”一注,即n\textsuperscript{o}~47时,我才被引得想起)。相反,我却没有觉察:在那篇谈我的短文里,“上同调”一词不曾出现,“概形”一词也一样。在我当时那种心不在焉的状态里,没有任何东西使我疑心:这篇平淡无奇、堆满夸张形容词的文本,竟充当着一份葬礼悼词的角色——而且是(外加)“以完美的手法呈献的”!手法完美到这种地步,以致我怀疑:这本小册子的读者们(它读来有点无聊,因为满纸都是场合所要求的、四面八方的蓄意涂脂抹粉,想必如此……)是否有人比我更早看出这一点——在我第一遍、第二遍阅读的时候。

这立刻汇入一条我不断重温的确认:每当我由于这样或那样的缘由,得以以一种\dmark{}\origpage{436}多少强烈而持续的注意,去端详某样我先前只满足于以“惯常的”、例行的注意“顺带”看过的事物——那些日复一日在我生活中掠过的大事小事——就会出现这样的局面。这种局面在禅观时期频频出现;禅观屡屡(而且多半“顺藤摸瓜”、并无蓄意)使我得以以前所未有的仔细,去检视某些白日或黑夜的事件(梦也在内):它们在我惯常的注意状态里或多或少未被察觉,或者其意味(往往清楚而分明)起初全然逃脱了我有意识的注意。

当我在这里说“多少强烈而持续的注意”,我心底所指的,其实是一双\textbf{醒着的眼睛},一双新的眼睛——不被思维的积习压重、也不被充当门面的“所知”压重的眼睛。只要由于这样或那样的缘由,我们得以用醒着的、专注的眼睛去看事物,它们便仿佛在我们眼前变了。在我们日常“注意”所呈递的事物那呆板光滑的表面的貌似平庸背后,我们骤然看见一道未曾料到的\textbf{深处}敞开、活跃起来。事物的这份深处的生活,不曾等我们费心去认识它才在——它无时不在;它是它们内在本性的一部分,无论所论的是数学对象、园子里的一片草坪,还是某个时刻在某人身上起作用的全部心灵力量。

\textbf{思维}是揭破并容我们探测表面之下的深处、事物的隐秘生活的一件工具——说它“隐秘”,只因为我们太懒,不肯看;太受抑制,不肯看见。它是一件有其长处、也有其短处与限度的工具。但无论如何,思维很少被当作发现的工具来使用。它最常见的功能,不是发现在我们身上、在事物之中的隐秘生活,而恰恰是遮掩它、把它凝固化。它是一件多用途工具,同时供“劳作的孩童”与“老板”使用。在这一位手里,它化作一面帆,善于捕捉我们欲望的风,把我们送往未知深处;在另一位手里,它凝成一只纹丝不动的锚,任漩涡与风暴都撼它不动……

反思本来有点走岔了;瞧,它回到了一个出发点——也就是昨天我停留过的那条确认:积习与烙印何等根深蒂固,使我远未发挥\dmark{}\origpage{437}自己的能力!(在这件事上,我还发现自己同伴有的是……)正是借着对葬礼的逐步发现——从LN~900那部文集这样扎眼的事实出发\fnote{8}{(*)}{见“一个梦的回忆——或Motif的诞生”一注,n\textsuperscript{o}~51,及随后的“下葬——或新父亲”一注。}——一份懒惰的注意终于醒了过来。读“拒受遗产——或矛盾的代价”一注(n\textsuperscript{o}~47),使我在5月12日第三遍(!)重读那两篇大名鼎鼎的“短文”。这一回,我总算注意到一个有点出格的细节:在那本庆典小册子里献给我的那篇酒神颂风格的小文里,通篇绝口不提“上同调”(也不提代数簇或概形)!这事足够滑稽,值得一脚注;我当即写了起来。写着写着,我又发觉一两处别的“滑稽”细节,它们还不曾咬住我的注意:那虽是第三遍阅读,它照旧停留于表面、机械——差不了多少,我不过是\textbf{重复}、\textbf{复制}此前的阅读而已。只是在写那本该是一条脚注、后来变成“葬礼悼词(1)”一注的文字时,我才渐渐被咬住了钩;一份\textbf{好奇}醒了,使我再一次回到那两篇文本上,这一回看得稍近了些。直到那一刻,方才所说的那种转变才发生:一道“深处”敞开了——在一场重大场合的锦衣华盖之下端出的酒神颂的平板门面背后,有一种炽烈的生命!正是这份好奇,把一道机械的、重复的、心不在焉的目光,变成一道“醒着”的目光……

而且,所说的“觉醒”也非一蹴而就;它是渐进的,随着那条(真不愧是)脚注里继续的反思的进程而成。老实说,直到那条注画上句点,它还不曾完成——尽管(我记得)当时夜已深,催我“收尾”\fnote{9}{(**)}{这一点想必尤其如此:就在当天,我已先经历了那场漫长而充实的“屠杀”反思(n\textsuperscript{o}~87);而且紧随其后的“葬礼悼词——或恭维”一注,在其结尾处还指涉了它。}。可我前脚刚画上那个点——或者至少到次日——我就发觉:葬礼悼词这个题目,我还远没有穷尽。\dmark{}\origpage{438}直到那一刻,我才\textbf{充分}感到:这两篇看似如此短小、如此平淡的文本,含义何等丰饶;干脆说,是两座真正的矿!而且,只要我肯倾听,它们要说的话,我还远没有巡完……

(9月25日)昨夜,反思又不得不中途收住——虽说它在我感觉里才刚刚起步。然而我已连续在打字机前坐了三个半钟头;一些隐约的小征兆开始向我示意:该起身动动了。

我清楚地记得头一回:我得以把一份“强烈而持续的注意”投向成文的文本,并且日复一日、一连数月,亲历一场“表面”的惊人蜕变——呆板而平板的表面活了起来,显出丰富而精确的意味,一道未曾料到的“深处”。那同时也是我头一场长程的禅观:怀着一场未知之旅的精神,旅程要多久就多久……出发的材料,是我父亲(流亡巴黎)与我母亲(其时还在柏林,带着当时五岁的我)之间1933/34年间的浩繁通信。我的目的,是“认识”我的父母。前一年我发现:我一生对他们怀有的那份敬佩——它终于凝成一种孝顺的虔敬——掩盖并维持着对他们的一份极大的无知。这份惊人的无知——我一生乐于把自己保持在其中——还是在那之后一年、1979年8月至1980年3月的那场长程禅观里,才向我显出其全部维度。

我于1979年7月间开始“整地”:其中主要是对这宗通信做第一遍通读,在我当时为一部“我自撰的诗作”\fnote{10}{(*)}{这部著作、以及它所代表的我生平的那段情节,见“非古鲁的古鲁,或三足马”一节(n\textsuperscript{o}~45)的结尾,及它所指涉的注n\textsuperscript{o}~43.}收尾之余进行。每天晚上,我花几个小时读三四封回信;兴趣当然是有的,而且当时我会毫不犹豫地说:读得很专注。然而我隐约发觉:我依旧陌生、外在于我所读的东西——真正的意味从我这里滑脱。我所读的常常相当离奇,\dmark{}\origpage{439}仿佛我眼前起居、招摇的这一男一女,与我自以为认识的那些人毫无瓜葛——我的记忆向我返还一幅清晰、分明、不可动摇的图像的那些人。由于缺乏一项耐心的、一丝不苟的、苛求的工作——随着前进逐步展开的——我只有震惊而已,为这些书信里(相对)不多的几处“够分量”咬住我表面注意的东西。如此登记下来的东西,只是叠加在“众所周知”之上;而那个“众所周知”,自我幼年起、直到那些天,始终是(当然,我从不曾察觉)我生活的、我的身份感的无形而不可动摇的基石。假定我当时就止步于这第一遍阅读,那么薄薄一层未消化的新“事实”——它就这样叠加在主宰层之上——想必很快会被磨蚀、冲走,在其后的岁月里不留多少痕迹。

在那项预备性工作之际,我的主要投资在别处——在一部吞噬了我大部分精力的著作的写作里。我分明意识到,在一项别的工作的边缘上做的工作有其限度;我迟早得从头到尾重新来过,做一项全力投入、逐件核查原始材料的工作。我预计那将是几个星期的事——实际上,我在其中连着泡了七个月:细细检读我父母留下的书信与文字;其中最“灼热”的部分,无疑是1933/34年的通信。而且,七个月之末,我终究还是收住了,发觉这个题目(“认识我的父母”)几乎取之不尽。如今更紧迫的,是\textbf{认识我自己}——借着我刚刚就我父母学到的这一切,并由此,至少间接地……认识我自己那段被遗忘的童年……

我刚花了近两个钟头,翻阅这场关于我父母的禅观——它始于1979年8月3日——的最初的笔记。与我匆匆记起的相反,我当时还没有认识到——除非极其模糊地——有必要“从头到尾”(如我方才所写)重看我此前一个月里读过的父母的书信与其他笔迹。至少,我在笔记里没有作这方面的任何暗示。在一两天的一场总结性反思——盘点这次阅读在我身上激起的多重印象,有些还相当混乱——之后,我全然没有摆出要以逐件核查原始材料的细致工作重做它的架势。\origpage{440}我宁可(仿佛理所当然地)接着以同样快的速度读\textbf{另一些}信(尤其是我父母1937/39年间的一宗浩繁通信),并继续一场由阅读印象喂养的平行反思。正是在8月与次月的进程中,顺藤摸瓜地,我开始懂得:对一封信(或一份生活的其他笔迹见证)下一番\textbf{功夫}是什么意思——那功夫使人得以把握其真正的、有时昭然若揭的意味;而那意味,写信的人却往往乐于无视、瞒过自己如同瞒过别人,神不知鬼不觉!——同时却又把它“摊在字里行间”,有时张扬、有时尖刻。而且,想必很少有这样的情形:那份影射或挑衅(有时凶狠……)竟不曾抵达收信人,不曾被他在某个层面上知觉并“照单收下”——尽管他自己也小心不让这份知觉、这份认识进入他目光的视野,而且他同样扬帆满舵,入的正是这同一出“神不知,鬼不觉!”的游戏。恰恰是那些最晦暗的段落——无一例外地,那些看似近乎低能(或疯狂……)、挑战一切合理解释的段落——在一双好奇的眼睛面前,显出意味最为丰饶:真正的矿,供给无可替代的钥匙,使人得以更深地穿入那一大堆表面的无意义背后那简单而分明的意味。这样的段落,在我父母的通信里、尤其在领舞的我母亲的信里,比比皆是;头一遍阅读——7月间——它们自然全从我头顶掠过。次月里,我开始此处彼处咬住它们。直到9月,几路旁证才使我明白:分明,在我本该从1933/34年的信里学到的东西中,我也许漏掉了某些本质的东西;它们把我带回那些信,催我对其中某些做第一遍“深读”。这一读,立刻把我自童年以来关于我父母其人、关于他们与我和我姐姐的关系曾经如何的图像,翻了个底朝天。

\section{书写的颂歌}

\notehead{102}{}

(9月26日)两天了,我一头扎进“自传式的追忆”——而我本来是要(趁热情已经冷却)续写某条关于某篇葬礼悼词的注的。我不知道这场枝蔓是否多少焐热了我的兴致!至少,该是到达我前天动笔时所瞄准的那个点的时候了——大方向是:“论阅读一份装作不说它要说之话的讯息的艺术”。\dmark{}\origpage{441}这类文本-讯息,远比我当年料想的常见……

不用说,只要人们还甘愿(一如我一生的大部分时光)把你听到或读到的一切当现款收下、按字面领受,并且无论在什么事、什么人身上,都不去寻找当事人明白说出的意图之外的别的意图,那么这门“艺术”的“如何”就无从谈起。相反,当人们有那种难以名状的感觉——某项声明、滔滔一番或一段叙述里,有什么“不对劲”,里面有猫腻,有什么东西“溜过去了”,而它照理是不该被说出的(您还想在这儿猜什么!)——问题就来了。有时,那是一种基本而令人错愕的知觉:一处不自洽、一件荒谬,有时如此巨大、又同时看似无从把握,以致仿佛挑战一切表述,濒临低能或谵妄的边缘。这类局面常常被焦虑压得满满——而每逢这样的局面,荒谬骤然闯入我的生活:一桩不可容忍、不可理解的、暗藏威胁的荒谬,每一次都撼动我那关于世界与我自己的宁静视野的根基——我对它的反应,千篇一律,是一股即刻涌来的焦虑(从不被承认为焦虑,却立刻被搅浑、藏掖在一场狂暴而绝望的愤怒的浪涛之下)。至少,直到我发现“禅观”之时是如此:自那以后,一份无畏而进取的好奇,拆除了这些愤怒与焦虑之浪的引信,并接替了它们……

正是好奇——即求知的欲望——使我在需要的压力下,自发地找到了这门破译一份混沌的见证文本的“艺术”——或者说得更谦逊些:找到了一套适合我那有限手段与我那份滞重的方法。尽管我又好奇又努力,这些含义深重的信,第一遍(甚至第二遍)读下来,本质的东西全从我头顶掠过——“我什么也没看出来”。有时,就某几处格外晦暗而令人错愕的段落,对着或许纷乱的印象写几句评语,笔锋所至,竟得以更深地进入一篇看似密不透风的文本的意味。写着写着,我有时不得不为了引用,抄录一些或长或短的段落:它们或以晦暗见称,或因为粗看之下出于这种或那种缘由给我“重要”的印象。\origpage{442}日子与星期流逝,我发觉:单是\textbf{全文照录}我所端详的文本的某一段,就以惊人的方式改变了我与该段的关系——朝着领会其真正意味的方向敞开。

这完全出乎意料,因为我最初的动机(至少在有意识的层面)纯粹是为了省事。我甚至记得,长年以来,我心里有一种压抑的不耐:把宝贵的光阴花在仅仅充当一名抄写员上!我按捺着焦躁,尽可能快地写,只盼早点抄完……可是,眼睛扫读成行的文字之快,与手逐字逐句转录它们之慢,没有共同的尺度。写得再快,“时间因子”也全然不同。而我疑心,这个“时间因子”并非以纯然机械的、定量的方式起作用——或者不如说,它只是一份更精微、更丰富的实在的一个侧面。至少在我这里,眼睛扫过别人构思并写下的文字的行动,与手一个字母一个字母、一个词一个词重写这同样一些文字的行为,也没有共同的尺度。可以肯定:手,与精神或思维之间,有一种深层的共生;随着书写之手的节奏,而且不带任何蓄意,精神不由自主地重新构成、重新思考同样的词——它们组接成承载着意味的句子,而这些句子又连缀成话语。只要一份求知的欲望鼓舞着这只誊抄字母、词语与句子的手,并同时鼓舞着那个以同一音高、在另一个层面上也“誊抄”着它们的精神——可以肯定,这双重行动,便在我这个甘当抄写者-执笔人的人与这份讯息之间,造成一种远比那件主要消极、不留任何有形踪迹的阅读之眼的行为亲密得多的接触。

这份摸索而来的直觉,与一条由来已久的确认同向——那就是:在我身上,工作者的思维的节奏(无论数学工作还是任何别的工作,包括我称为“禅观”的工作),最常见(倘若不是永远)是书写之手的节奏,而绝不是阅读之眼的节奏\fnote{11}{(*)}{这一情势在我身上起作用,似比在我的大多数数学家同行身上更强;它当年使我难以融入Bourbaki小组的集体工作:我发现自己全然无力跟上那些以它们进行的速度宣读的报告。况且,我从不曾真正喜爱过\textbf{读}数学文本,哪怕再漂亮。我理解数学的自发方式,始终是\textbf{做}它们,或\textbf{重做}它们(必要时,此处彼处借同事们供给的想法与提示,退而求其次,借书……)。}。而我的手(或有时,由我的双手操纵的打字机……)\dmark{}\origpage{443}留下的\textbf{笔迹}——它以那份不慌不忙、却也从不磨蹭的思维前进的节奏写出——是这份思维不可或缺的物质支撑:既是它的“声音”,也是它的“记忆”。况且我疑心,在大多数(倘若不是全体)“智力工作者”身上,情形或多或少也是如此(也许程度稍轻)。

\section{孩子与海——或信与疑}

\notehead{103}{}

(9月27日)不管怎样,事实摆在这里:一如我唯有经由书写才能“进入”一门数学理论,我也几乎唯有经由\textbf{重写},才开始进入一份文本-讯息、一份讯息的“字里行间”。我的第一场“就文本而做”的禅观工作发生转变——一道表面的平庸开始向一道活的深处敞开,荒谬开始取得意味——恰恰\textbf{从那一刻起}:我开始全文重写那份讯息,或(当其篇幅大到无法处理时)重写那些嗅觉使我感到关键的段落。

人们会对我说:既然没有可靠的“客观”判据来担保一份“诠释”的有效性——它被呈献为对(比如说)一份文本的一项(所谓的?)“工作”的结果或终点——人们尽可以让任何文本、任何话语说出自己想说的任何东西,给它编造自己喜欢的随便什么“讯息”。这当然千真万确,例子肯定比比皆是!况且我怀疑(或许除了一门划界分明如史学的学科——即便如此……),这样的判据是否能提炼出来。反正这也无济于事:既不能阻止任何人随意编造异想天开的诠释,也不能使任何人得以探测并发现一份讯息、一个局面、一桩事件的真正意味。规则与判据是一门\textbf{方法}的配料;方法自有其用处与重要性(不过常被高估,以牺牲一些性质全然不同的其他因子与力量为代价)——作为科学或技术知识之发展中的发现与巩固的工具,也作为任何一门手艺(开车或修车,等等)中的工具。相反,在认识自己与他人的层面上,方法的角色变得全然次要:“总管事务”自会跟上来,当本质的东西在场时。而取法于一项方法、或从它出发,乃至死死抱住它,丝毫无助于那份更本质的东西的显现——恰恰相反!

换个说法:那个出发去寻找某样\textbf{预先决定}\dmark{}\origpage{444}的东西的人(他将称之为“真的”,或“真理”),将毫不费力地找到它,甚至证明它,直到自己全然满意;而且他肯定会在路上遇到这一个或那一个同道,甚至一整群,欢欢喜喜与他结盟、分享信念与满足。他像那位捕蝶人:出发时网里就装着一只漂亮的蝴蝶(说不定还是制成标本的),打猎归来,喜滋滋地(而且全然满意地)把它从网里掏出来。

还有另一种人:面对一个未知,如一个赤身的孩子面对大海。当孩子想认识它,他便走进去,便认识了它——无论它温凉,平静或起伏。被某个未知之物吸引、出发去认识它的人,肯定会多少认识它。带网也罢,不带网也罢,他会找到真,或无论如何,找到\textbf{一点}真。他的错误一如他的发现,都是他行程中的台阶;或者不如说,是他与所欲认识者的\textbf{恋情}中的台阶。

我深知自己在说什么,因为我这一生,轮流地,既当过那位捕蝶人,也当过那个赤身的孩子。把二者区分开来,毫无困难。我怀疑“客观判据”在这里能帮上什么忙——事情比那简单得多!只消用自己的眼睛……

同样毫无困难的,是分辨我刚谈到的这段行程中相继的台阶、相继的澄清阶段:从那一“死”的阶段——没有任何预感浮上意识,使人无从疑心“某种东西”,超出困倦的眼睛呈递给我们的那层平面而无定形的表面——直到经由相继的“觉醒”,抵达对这份“某种东西”越来越精微、越来越贴身、越来越完整的把握。无论所论的是在数学事物之发现中的行程,还是在对自己与他人的发现中的行程,其本性并无本质不同。在一份\textbf{认识}中的\textbf{前进}之感——它渐渐加深(哪怕穿过一份错误之积累,耐心地、不倦地不断更正)——在这后一种情形,一如在前一种情形,同样无可辩驳。

这份\textbf{安然的确信}——这是一份内在心性的一面;其另一面,是一份\textbf{向怀疑敞开}:一种好奇的态度,排除对自己的错误的任何恐惧,使人得以不断侦察并更正它们。这一双重立足——这份接纳怀疑与从事发现同样不可或缺的\textbf{信}——的本质条件,是对这项探索将“产出”什么毫无恐惧(无论明显的\dmark{}\origpage{445}还是暗藏的):尤其是,毫无这样的恐惧——怕我们即将发现的实在动摇我们的确信或信念,怕它使我们的希望幻灭。这样的恐惧,如同一道深深的瘫痪,瘫痪我们的创造官能、我们更新的力量。我们可以在艰辛与痛苦中发现并更新自己;却不能在恐惧中发现——恐惧那即将被认识者、即将诞生者。(正如一个人若在某个时刻惧怕一个女人、或惧怕把他带入她体内的那件行为,便不能认识她、使她受孕。)这样的恐惧,在一项科学研究、或任何其主题并不多少深刻牵涉我们自身的研究的语境里,想必相对罕见。相反,当事关对自己或对他人的发现时,它却是最大的绊脚石。

然而,伴随着一次发现——无论大小——的那份感觉,在对自己的发现或对他人的发现的情形里,一如在一场非人称的研究(比如数学)的语境里,同样无可辩驳。我已有机会影射过这份感觉。它是刚刚发生的某事的知觉在情绪层面的反照——某样\textbf{新}东西的显现——而这份“某样东西”,显得与比如说一个数学陈述、或一个此前从未有人想过的观念或证明的显现,同样具体可触、同样无可辩驳(请恕我重复!)。况且我觉得,难以把伴随某次特定发现的那份感觉,与我方才谈到的伴随整场研究的前进之感,区分或分割开来。大的与小的发现,如同把一份前进落实下来的相继\textbf{台阶},如同我们必须相继越过的\textbf{门槛}。\origpage{446}前进不是别的,正是这一连串门槛的跨越、每一台阶向下一台阶的升进。

那份“感觉”——或更好地说,那份映照、还原这一过程的知觉——是一个可靠、不容置疑的“判据”:我不记得它曾把我引入歧途,无论在数学还是在禅观;我不记得曾带着事后的认识确认:这份感觉曾是幻象。它常常使人毫无残余地分辨真与假,或辨出假中之真、与那份自命为真者之中的假。但它尤其是一切真正探索中不可替代的\textbf{向导}——一位随时准备(只要我们肯费心求教)告知我们是否走错了路、抑或正走在好路上的向导。

对这份可靠向导的倾听状态,依我看,不是别的,正是我在反思的另一处\fnote{12}{(*)}{见“严与苛”一节(n\textsuperscript{o}~26),我在其中把“严格”谈作一场探索中“对每一时刻当下\textbf{理解之品质}的精微留意”。}所命名的“严格”。这份严格,依我看,无论所论的是一场数学探索中的苛求,还是自我认识中的苛求,并无本质不同——没有后者,那样的认识便无从谈起。但不用说,这绝不意味着:这份严格在某一智力工作层面的在场,就担保或标志着它在自我与他人的认识中的在场。实际上,相反的情形才是真的;我从无数机缘里——从我自身开始——确认了这一点。在那个领域里,我这里所说的“严格”,是与禅观同时来到我的生活中的。或者不如说,我实在无从区分二者。我生活中的禅观时刻,不是别的,正是我以这样一份对自己极端苛求的心境检视自己的时刻——多半经由我与他人的关系。
